\documentclass[11pt,a4paper]{article}
\usepackage[utf8]{inputenc}
\usepackage[T1]{fontenc}
\usepackage[margin=2.2cm]{geometry}
\usepackage[dvipsnames]{xcolor}
\usepackage{titlesec}
\usepackage{enumitem}
\usepackage[hidelinks]{hyperref}
\definecolor{accent}{rgb}{0.11,0.26,0.49}
\input{glyphtounicode}\pdfgentounicode=1
\titleformat{\section}{\large\bfseries\color{accent}}{}{0em}{}
\titlespacing*{\section}{0pt}{1.2em}{0.4em}
\setlist[itemize]{leftmargin=1.4em,itemsep=2pt,topsep=2pt}
\setlength{\parindent}{0pt}
\setlength{\parskip}{0.4em}

\begin{document}

\begin{center}
  {\LARGE\bfseries\color{accent} Pour la France, la contrainte externe n'a pas disparu}\\[0.3em]
  {\large André Cartapanis (2022)}
\end{center}

\vspace{0.4em}
{\small\itshape Note publiée par le Cercle des économistes, en prélude aux 22\textsuperscript{e}
Rencontres Économiques d'Aix-en-Provence. Auteur : André Cartapanis, professeur à
Sciences Po Aix (p.~3).}

\section{Résumé exécutif}
On affirme souvent qu'en union monétaire la « contrainte externe » a disparu.
L'auteur soutient que c'est faux pour la France : son déficit extérieur, lié à ses
déficits budgétaires, pèse sur la soutenabilité de sa dette vis-à-vis du reste du
monde et sur ses conditions de financement (p.~3). À l'appui de données du FMI, il
montre que la France cumule un déficit courant récurrent, un déficit budgétaire
persistant et un endettement externe croissant, ce qui réduit ses marges de
manœuvre — sans pour autant conclure à une crise imminente (p.~15).

\section{Contexte \& problématique}
Le déficit commercial français de mars 2022 (plus de 12 milliards d'euros, soit
environ 100 milliards en données glissantes sur douze mois) relance l'attention sur
la balance des paiements (p.~4). La question centrale : l'appartenance à la zone
euro protège-t-elle réellement la France de la contrainte externe, ou celle-ci
s'est-elle simplement déplacée vers l'endettement et son financement (p.~3, p.~12) ?

\section{Points clés}
\begin{itemize}
  \item La France enregistre \textbf{structurellement} des déficits de balance des
  paiements courants depuis 2014, alors que la zone euro dégage des excédents
  systématiques — un « paradoxe » pour un pays membre de cette zone (p.~4, p.~6).
  \item D'un point de vue macroéconomique, un déficit courant traduit une épargne
  interne insuffisante pour financer l'investissement et le déficit budgétaire,
  d'où un recours à l'endettement externe (p.~7).
  \item Un déficit courant n'est pas « bon » ou « mauvais » en soi : tout dépend de
  ses causes et de sa soutenabilité dans le temps (p.~8-9).
  \item Le principal risque associé à un déficit courant est le \textit{sudden stop} :
  un arrêt brutal des entrées de capitaux si la perception du risque change (p.~9-10).
  \item La crise de la zone euro (2010-2014) prouve que l'appartenance à l'union
  monétaire n'a pas supprimé les crises de balance des paiements : c'était une
  « illusion » des fondateurs de l'euro (p.~10-12).
  \item Pour la France, les déficits courants sont « indiscutablement liés aux
  déficits budgétaires » qui pèsent sur l'équilibre épargne-investissement (p.~13).
  \item La France fait partie des rares économies à cumuler déficit courant
  récurrent, déficit budgétaire non résorbé et fort endettement externe, le tout
  avec une croissance potentielle faible (p.~15).
\end{itemize}

\section{Données \& faits marquants}
\begin{itemize}
  \item Déficit commercial de mars 2022 : > 12 milliards d'euros ; ~100 milliards
  en cumul sur douze mois (p.~4).
  \item Déficit courant 2022 projeté par le FMI : ~54 milliards de dollars, soit
  ~1,8\,\% du PIB (p.~4-5).
  \item Déficit courant 2020 (Banque de France) : 43,7 milliards d'euros (1,9\,\% du
  PIB), plus haut niveau depuis 1982 (p.~6).
  \item Excédents courants de la zone euro : entre 250 et 400 milliards de dollars
  (1,9\,\% à 3,2\,\% du PIB agrégé) (p.~6).
  \item Déficits courants 2022 : États-Unis 877,8 milliards \$ (3,5\,\% du PIB) ;
  Royaume-Uni 184,4 milliards (5,5\,\% du PIB) (p.~6).
  \item Déficit budgétaire français : ~3-4\,\% du PIB depuis 2013, puis 9,1\,\% en
  2020 et 7\,\% en 2021 (p.~13).
  \item Endettement public français : 83\,\% du PIB en 2013, 88,8\,\% en 2019,
  au-delà de 100\,\% après la pandémie (p.~13-14).
  \item Endettement net externe : 547 milliards d'euros en 2017 → 694 milliards en
  2020 (de 21,1\,\% à 26,4\,\% du PIB) (p.~14).
  \item La position extérieure nette dépasse le seuil de 30\,\%, sans atteindre le
  seuil officiel de 35\,\% qui déclenche la procédure européenne de déséquilibre
  macroéconomique excessif (p.~14-15).
\end{itemize}

\section{Termes \& concepts clés}
\begin{itemize}
  \item \textbf{Contrainte externe} : limite imposée à un pays par la nécessité
  d'équilibrer ses échanges avec le reste du monde (sa balance des paiements).
  \item \textbf{Balance des paiements courants} : solde des échanges de biens,
  services et revenus d'un pays avec l'étranger.
  \item \textbf{Déficit / excédent courant} : un déficit signifie qu'un pays dépense
  plus qu'il ne produit et doit s'endetter auprès de l'étranger ; un excédent,
  l'inverse (p.~7).
  \item \textbf{Sudden stop} : interruption brutale des entrées de capitaux
  étrangers, qui peut déclencher une crise (p.~9-10).
  \item \textbf{Position extérieure nette} : différence entre ce qu'un pays détient
  à l'étranger et ce qu'il doit à l'étranger (p.~14).
  \item \textbf{TARGET2} : système de règlement entre banques centrales de la zone
  euro, révélateur des tensions de balance des paiements internes à la zone (p.~12).
  \item \textbf{Global saving glut} : « excès d'épargne mondiale », concept de Ben
  Bernanke évoqué pour le risque futur d'une raréfaction de l'épargne disponible
  (p.~16).
\end{itemize}

\section{À retenir}
L'auteur ne formule pas de recommandations chiffrées : en union monétaire, la
dévaluation n'est plus disponible, et le problème français ne vient pas d'un
déficit de compétitivité-prix mais d'un déséquilibre macroéconomique d'ensemble,
d'une croissance potentielle faible et d'une spécialisation insuffisante dans les
produits à haute valeur ajoutée (p.~17). Sa conclusion : la France ne peut pas
reporter indéfiniment l'ajustement de ses comptes publics et de sa balance des
paiements ; la contrainte externe n'a pas disparu (p.~17).

\vspace{0.3em}
\textbf{Limites / portée :} il s'agit d'une note d'analyse argumentée (point de vue
d'auteur), pas d'une étude neutre. L'auteur souligne lui-même que la théorie
économique « est incapable » de dire s'il existe un seuil de déclenchement d'un
\textit{sudden stop}, et que l'histoire « ne se répète jamais à l'identique » (p.~15).

\end{document}
